\scanpage{185}
\uline{Groupe de Picard des cones} \uline{\unsure{épointés|épointis}}

$X\subset P_k^{r}$, cône projetant (affine) $C$, $a$ l'\deleted{origine} sommet, $C_a=\operatorname{Spec}\underline O_{C,a}$.

$C-a$ est un fibré principal sur $X$, de groupe $G_m$, dont la classe est celle $\xi\in\operatorname{Pic}(X)$ de $\underline O_X(1)$. Donc \uline{si} $X$ \uline{est régulier}
\[
\operatorname{Pic}(C-a)\xleftarrow{\sim}\operatorname{Pic}(X)/\mathbb Z.\xi
\]

Considérons le morphisme naturel
\[
C_a-a\longrightarrow C-a,
\]
il induit un hom canonique
\diagram{\begin{tikzcd}[ampersand replacement=\&,row sep=tiny]
\operatorname{Pic}(C-a)\arrow[r]\&\operatorname{Pic}(C_a-a)\\
\&{\displaystyle\lim_{U\ni a}\operatorname{Pic}(U-\unsure{a})}\arrow[u,phantom,"\rotatebox{90}{$\simeq$}" description]
\end{tikzcd}}

Je dis que c'\uline{est un isomorphisme si $X$ est régulier}. Comme pour $U=C-Y$ voisinage de $a$, on a
\[
\operatorname{Pic}(U-a)\simeq\operatorname{Pic}(C-a)/\operatorname{Im}\mathbb Z^{(I)}
\]
où $I$ est l'ens des comp. irréd. de $Y$,
% archive sheet number: 184
